% ================================================================
% ==== FILE: content/cycles/cycles_k4.tex
% ================================================================

\subsubsection{Zyklen auf \texorpdfstring{$K_4$}{K_4}}
\label{sec:cycles-k4}

Ebene$K_4$entspricht dem protozellulären Kontinuum:
eine chemisch selbstorganisierte Domäne, die von einem semipermeablen umgeben ist
membrane $\partial\Omega(K_4)$.
Zyklen definieren die persistente dynamische Struktur der Protozelle:
Aufrechterhaltung von Gradienten, Stoffwechselumsatz, Membranreparatur,
Energieumwandlung, Replikationsgenauigkeit und frühe Protoerregbarkeit.
Sie bestimmen die Lebensfähigkeit der Protozelle und die Möglichkeit
des Übergangs zum Erregbarkeitsniveau$K_5$.

\subsubsection{Klassen von Zyklen auf \texorpdfstring{$K_4$}{K_4}}

Nach der Formalisierung Biologie U0.3b,$K_4$unterstützt sechs
grundlegende Familien von Zyklen.

\subsubsection{1. Membranwartungszyklen$C_{\mathrm{mem}}$}

Diese Zyklen reparieren, stabilisieren und formen die Membran neu.
Dazu gehören:
\[
  J_{\mathrm{lipid}},\quad
  J_{\mathrm{in/out}},\quad
  P_{\mathrm{mem}},\quad
  \gamma_{\mathrm{edge}},\quad
  \kappa_{\mathrm{bend}},
\]
where $\gamma_{\mathrm{edge}}$ is the line tension of membrane pores
and $\kappa$ is the bending modulus (the corresponding derivation note).

Die Zyklusstruktur ist:
\[
  \partial\Omega \to
  \mathrm{stretching} \to
  \mathrm{repair} \to
  \partial\Omega.
\]

Stabilität erfordert:
\[
  T_{\mathrm{mem}}(t) < \Theta_{\mathrm{mem}},
\]
where $T_{\mathrm{mem}}$ measures osmotically and mechanically
induzierte Spannung.

\subsubsection{2. Gradientenerhaltungszyklen$C_{\mathrm{grad}}$}

Die Membran führt neue Achsen ein:
\[
  A_{\mathrm{in/out}},\;
  A_{\mathrm{perm}},\;
  A_{\mathrm{grad}}.
\]
Dementsprechend regulieren Gradientenzyklen:
\[
  P_{\mathrm{grad}}=
    (\Delta \mathrm{pH}, \Delta \mathrm{ion}, \Delta \mathrm{redox}),
\]
mit Flüssen:
\[
  J_{\mathrm{pump}},\;
  J_{\mathrm{leak}},\;
  J_{\mathrm{channel}}.
\]

Der Zyklus hat die Form:
\[
  \Delta G \to \mathrm{pump} \to
  \mathrm{leak} \to \mathrm{dissipation} \to
  \Delta G,
\]
where $P_{\mathrm{grad}}$ must remain in the viable domain:
\[
  |P_{\mathrm{grad}}| < \Theta_{\mathrm{grad}}.
\]

Diese Zyklen sind Vorläufer der elektrochemischen Struktur von$K_5$.

\subsubsection{3. Energie- und Redoxzyklen$C_{\mathrm{energy}}$}

From the corresponding derivation note, the energetic architecture includes:
\[
  P_{\mathrm{energy}},
  P_{\mathrm{redox}},
  P_{\mathrm{grad-energy}},
\]
und fließt:
\[
  J_{\mathrm{energy}},\;
  J_{\mathrm{redox}},\;
  J_{\mathrm{grad}}.
\]

Der protometabolische Zyklus ist:
\[
  \Delta G_{\mathrm{in}} \to
  \mathrm{activation} \to
  \mathrm{transfer} \to
  \mathrm{dissipation} \to
  \Delta G_{\mathrm{in}}.
\]

Existenz erfordert:
\[
  T_{\mathrm{cycle-energy}} < \Theta_{\mathrm{cycle-energy}}.
\]

Diese Zyklen stabilisieren das innere Umfeld und unterstützen
Gradientenerhaltungszyklen.

\subsubsection{4. Stoffwechselumsatzzyklen$C_{\mathrm{metabolic}}$}

Aufbauend auf der RAF-Vorläuferstruktur von$K_3$, $K_4$Formen
lokalisierte protometabolische Schleifen, eingebettet in die Membran:
\[
  M_{\mathrm{in}} \xrightarrow{r_1}
  M_{\mathrm{act}} \xrightarrow{r_2}
  M_{\mathrm{out}} \xrightarrow{r_3}
  M_{\mathrm{in}}.
\]

Sie tragen dazu bei:
\[
  P_{\mathrm{in}},\; P_{\mathrm{chem}},\; P_{\mathrm{stability}},
\]
und brauchbare Konzentrationsbereiche aufrechtzuerhalten.

Störung:
\[
  T_{\mathrm{metabolic}}>\Theta_{\mathrm{stability}}
\quad\Rightarrow\quad
  \Omega(K_4)=\varnothing.
\]

\subsubsection{5. Informationsstabilitätszyklen$C_{\mathrm{info}}$}

From the replication accuracy block (the corresponding derivation note),
Informationstragende Polymere müssen Folgendes erfüllen:
\[
  P_{\mathrm{copy}}=(1-\varepsilon)^L \ge P_{\min}.
\]

Die Mutations-Korrektur-Selektionsschleife:
\[
  \mathrm{template}
  \to \mathrm{copy}
  \to \mathrm{mutation}
  \to \mathrm{correction}
  \to \mathrm{selection}
  \to \mathrm{template}
\]
bildet einen echten Kreislauf.

Seine Lebensfähigkeit entspricht:
\[
  \varepsilon < \Theta_{\mathrm{error}}.
\]

Oberhalb der Fehlerschwelle die Informationsstruktur der Protozelle
bricht zusammen.

\subsubsection{6. Proto-Erregbarkeitszyklen$C_{\mathrm{exc}}$}

Based on the corresponding derivation note, early channels and gradients produce
Zünd-Front-Zyklen („Proto-Aktionspotentiale“):
\[
  P_{\mathrm{grad}} \to
  \mathrm{local\ ignition} \to
  \mathrm{front\ propagation} \to
  \mathrm{recovery} \to
  P_{\mathrm{grad}}.
\]

Die Zündung erfordert:
\[
  G_{\mathrm{lat}} > \Theta^{\mathrm{crit}}_{\mathrm{lat}}
\quad\text{and}\quad
  C_{\mathrm{mem}}>C_{\mathrm{mem}}^{\mathrm{crit}}.
\]

These cycles define the soft transition $K_4\to K_5$
(Polnisch~W3-Speicher).

\subsubsection{Metriken von Zyklen bei \texorpdfstring{$K_4$}{K_4}}

\paragraph{Length.}
\[
  L(C)=\oint d\Omega,
\]
where $\Omega$ is the state space of internal composition.

\paragraph{Efficiency.}
\[
  C_{\mathrm{eff}}
    =\frac{\int J\cdot dA}{L(C)},
\]
mit Strömungen$J$über:
\[
  A_{\mathrm{mem}}, A_{\mathrm{grad}}, A_{\mathrm{metabolic}}.
\]

\paragraph{Stability.}
\[
SC)
   = \min_{t\in C}
       d\bigl(
         \Omega(K_4),\;
         \partial\Omega(K_4)
       \bigr).
\]

\paragraph{Weight.}
\[
  w(C)=F\bigl(
        C_{\mathrm{eff}},\;
        S(C),\;
        \mathrm{energy\ turnover},\;
        P_{\mathrm{grad}}
      \bigr).
\]

\subsubsection{Zeit- und Wiederholungsstruktur}

Die definierenden Zeitskalen von$K_4$Zu den Zyklen gehören:

\[
  \tau_{\mathrm{pump}},\quad
  \tau_{\mathrm{leak}},\quad
  \tau_{\mathrm{repair}},\quad
  \tau_{\mathrm{metabolic}},\quad
  \tau_{\mathrm{exc}}.
\]

Das Kontinuum besitzt Zeit, wenn diese der stabilen Wiederholung genügen:
\[
  \tau_i<\infty \quad\forall i.
\]

\subsubsection{Zusammenbruch und Lebensfähigkeit von \texorpdfstring{$C(K_4)$}{C(K_4)}}

The death of $K_4$ (the corresponding derivation note) occurs if any of:

\[
  \Theta_{\mathrm{mem}},\;
  \Theta_{\mathrm{grad}},\;
  \Theta_{\mathrm{cycle-energy}},\;
  \Theta_{\mathrm{error}},\;
  \Theta_{\mathrm{stability}}
\]
wird verletzt.

Alle Zyklen brechen gleichzeitig:
\[
  C_j \to \varnothing,
\quad
  \Omega(K_4)=\varnothing.
\]

\subsubsection{Übergang zu \texorpdfstring{$K_5$}{K_5}}

Cycles $C_{\mathrm{exc}}$ and $C_{\mathrm{grad}}$ gradually
versteifen sich zu temporal-elektrischen Schleifen, passend zu den
Proto-Spike-Struktur von$K_5$:
\[
  C_{\mathrm{exc}}\longrightarrow C_{\mathrm{spike}}.
\]

Dieser Übergang ist kontinuierlich (C$^1$), wie gewährleistet durch
die polnische ~W3-Soft-Threshold-Formulierung.
